% !TEX root = ../../mixed-sklyanin-cubics.tex
% Current result: 2026-11-22 / v58.
\subsection{固定した楕円因子の周期付き境界比較}
\label{ebk:subsec:boundary}

上で固定した実因子 $S_i$ とその点データ
$(E_i,L_i,t_{p_i})$ を用いる。
各 $f_i$ は同じ偏極曲線 $(E_i,L_i)$ の任意の非零方程式であり、
非可換側の $w_i=a_i\Omega_{i,0}$ と可換側の $a_if_i$ を
それぞれ実際の因子方程式として固定する。
以下の指定周期とは、これらの方程式を用いた
graded MF と Orlov window から移した三乗周期をいう。

\begin{lemma}[任意の固定楕円因子に対する実境界核]
\label{ebk:lem:boundary}
$Z=V(G)\subset\mathbb P^{N-1}$ は滑らかであり、
境界の残余線束圏 $\mathcal K_S$ は全 coherent 残余圏に一致する。
両者は smooth proper で、$\C$ 上に実際の dg Morita 核
\[
 E:\mathcal K_S\xrightarrow{\sim}\Ku(Z)
\]
と rotation の比較射
\[
 \eta:E\varrho_S\xrightarrow{\sim}\varrho_ZE
\]
が存在する。$\eta$ を三回合成した同型は、
$\varrho_S^3\simeq[2]$ と $\varrho_Z^3\simeq[2]$ の
上で指定した周期を保つように取れる。
また、ブロック別スカラーから来る
$Z_0=(\mu_3^k)/\Delta\mu_3$ の指定作用も coherent に保つ。

因子ごとの二捩れ、異なる $j$ 不変量、同種関係や CM を許し、
$3\mid k$ の場合も含む。
\end{lemma}

\begin{proof}
各 $f_i$ の三つの偏微分は、対応する affine 三次元空間で原点以外に
共通零点を持たない。全 $a_i$ が非零で変数ブロックが独立なので、
$G$ の affine critical locus も原点だけであり、$Z$ は滑らかである。

各 $S_i$ は二次 Koszul AS-regular 三次元代数であり、
実際の正則中心元 $\Omega_{i,0}$ による商は
$B(E_i,L_i,t_{p_i})$ である。
主稿の定理\paperref{main}{mbnp:thm:product}の有限 MF 外積と
系\paperref{main}{mbnp:cor:coherent}の飽和を適用できる。
その仮定はこれらの実因子の性質であり、
$j$ の相異性や $2p_i\ne0$ を要求しない。
具体的には、まず三つの augmentation shift が生成する残余線束圏について
有限外積と Shamash 分解により積表示を得る。
積の smooth proper 性と有限群の同変化から残余線束圏の
smooth proper 性を得た後に、同系の coherent 飽和を適用する。
従って全 coherent 圏の生成性を外積比較の前提にしていない。
可換側にも同じ外積と window を適用すると、$Y=\prod_iE_i$ に対し
\begin{equation}\label{ebk:eq:equivariant-products}
 \mathcal K_S^{Z_0}\simeq\Perf(Y)
       \simeq\Ku(Z)^{Z_0}
\end{equation}
という実 dg Morita 比較を得る。
可換側は主稿の系\paperref{main}{mbrp:cor:comparison-field}の比較である。

ここで $Z_0$ は指定されたブロック別スカラー作用である。
その作用は
\[
 \Gamma_k=\{(\lambda_i)\in\mathbf G_m^k:
                    \lambda_i^3=\lambda_j^3\text{ for all }i,j\}
\]
による graded MF と、共通 grading scalar の自明化から得られる。
この構成は実因子の全自己同値群の分類を使わない。
以下の証明は、任意の実因子を扱う主稿の
定理\paperref{main}{pps:thm:elliptic-existence}の証明と、
定理\paperref{main}{csc:thm:comparison}の同変双加群の構成を、
この境界に適用したものである。

右加群の規約では、因子ごとの内部 shift はそれぞれ
\[
 R_{i,p_i}=T_{\OO_{E_i}}\circ(t_{p_i})_*\circ(-\otimes L_i),
 \qquad
 R_{i,0}=T_{\OO_{E_i}}\circ(-\otimes L_i)
\]
である。主稿の楕円 window の補題
\paperref{main}{lem:nfm-twisted-elliptic-window}と
命題\paperref{main}{mbnp:prop:rotations}は、これらを
実方程式からの指定三乗周期とともに同定する。
$\C$ 上では $p_i$ を表す実際の次数零線束 $M_i$ が存在し、
主稿の定理\paperref{main}{thm:nfm-elliptic-normalized-intertwiners}により
\[
 F_i=-\otimes M_i,\qquad
 \eta_i:F_iR_{i,p_i}\xrightarrow{\sim}R_{i,0}F_i
\]
を取れる。
両辺は可逆な perfect Fourier--Mukai 核なので、その自己準同型環は
$\mathrm{HH}^0(E_i)=\C$ である。
$\eta_i$ の三回の合成と両側の指定周期との差は
$u_i\in\C^\times$ であり、$\eta_i$ を $v_i$ 倍すると
差は $v_i^3u_i$ になる。
$v_i^3=u_i^{-1}$ と選んで指定周期を合わせる。
これは主稿の補題\paperref{main}{mbrc:lem:factor-correction}の実構成であり、
自然な方程式直線比較の非消滅を仮定しない。
特に二捩れでその自然な比較が零になる場合にも、
ここで選んだ非零方程式 $f_i$ をそのまま使える。

基準因子 $b$ を一つ選ぶと、式\eqref{ebk:eq:equivariant-products}の
双対指標作用は
\[
 Q_{i,p}=R_{i,p_i}^{-1}R_{b,p_b},\qquad
 Q_{i,0}=R_{i,0}^{-1}R_{b,0}\qquad(i\ne b)
\]
で与えられる。
$F=\boxtimes_iF_i$ と上の $\eta_i$ は
$FQ_{i,p}\simeq Q_{i,0}F$ を定める。
異なる因子の交換は外積の交換則によって保たれ、
$Q_i^3\simeq[-2][2]\simeq\mathrm{id}$ の図式は
合わせた指定周期によって保たれる。
比較射の次数は零で、周期の shift は偶数なので、交換の符号も一致する。
双対群の元を $\prod_{i\ne b}Q_i^{d_i}$、$0\le d_i<3$ の正規形で書けば、
その積は交換と三乗の消去で定まる。
元の MF の中心周期の coherence によりこれらの操作の順序は結果を変えず、
$F$ は全群法則の coherence を intertwine する。
ここで使うのは全次数格子
\[
 \Lambda_k=\mathbb Z^k/\langle3e_i-3e_j\rangle
\]
であり、対称な rotation 分裂を選ばない。
従って $3\mid k$ の場合も同じ証明が成立する。

この導来同変核を実 dg 核として実現する際は、
元の $Z_0$-同変 graded MF 模型へ戻る。
両模型上の $\widehat Z_0$-作用は線型化の character twist による
厳密な dg 作用であり、その間の perfect 双加群の dg 圏にも
両側の twist による厳密な作用がある。
式\eqref{ebk:eq:equivariant-products}で Morita 移送した $F$ は
この双加群圏の対象で、上で確認した図式が
その導来圏での同変構造を与える。
従って\cite[Lemma~8.6]{MEDElagin}を
この\emph{既に厳密な dg 作用を持つ双加群圏}に適用できる。
有限群の誘導、平均化と冪等完備化により、同変構造を
実際の同変 perfect 双加群として取れる。
同補題だけから任意の三角圏作用の dg lift を導いているのではない。
忘却した核は Morita 同値なので、得られた同変核も同値を与える。
\cite[Theorem~4.2, Theorems~8.7 and~8.9]{MEDElagin}の
二重同変化とその dg 実現により、$Z_0$-作用を保つ
$E:\mathcal K_S\simeq\Ku(Z)$ を得る。

最後に、基準因子の $R_{b,p_b}$ と $R_{b,0}$ は、
元の rotation の線型化を積表示へ移したものである。
$FR_{b,p_b}\simeq R_{b,0}F$ の射は、同じ因子比較と外積交換から作られ、
各 $Q_i$ の同変図式と可換する。
従ってこの射も同変双加群圏の射として持ち上がり、
二重同変化後に $\eta$ を与える。
三回の合成の図式は既に合わせた周期の図式であるため、
指定周期も保たれる。
\end{proof}

\begin{remark}[後続の comonad 比較への適用範囲]
\label{ebk:rem:scope}
補題\ref{ebk:lem:boundary}は固定した複素数体値の実因子についての存在結果である。
一度選んだ $E$ は、これらの因子を固定して変形係数だけを動かす
任意の係数環へ定数族として底変更できる。
この $E$ による $K\mapsto EKE^{-1}$ は endofunctor の合成を保つ
Morita 移送なので、実随伴の counit、余乗法と全高次 coherence を運ぶ。
rotation 比較により $\varrho_S^2[-1]$ も
$\varrho_Z^2[-1]$ へ同定される。
これが後続の固定 counit の arity 計算に必要な境界入力であり、
Heisenberg 同変性や異なる因子間で共通な scalar 正規化は不要である。
因子自体を動かす楕円モジュライ上では、実 Picard 線束と周期の三乗根を
族として選ぶ問題が別に残る。
\end{remark}
